\chapter*{Introducción}
\addcontentsline{toc}{chapter}{Introducción}

La Universidad Distrital Francisco José de Caldas ha procurado que sus estudiantes movilicen, entre otras competencias,
el uso de la tecnología. Además, en el proyecto curricular de Ingeniería Eléctrica de la Facultad Tecnológica, algunos documentos institucionales señalan que se propone la implementación de un laboratorio de cómputo especializado en matemáticas, el cual contará con paquetes tales
como Matlab, Mathematica y \emph{software} libre.
También se busca que el estudiante, en su proceso de aprendizaje autónomo, refuerce los conocimientos vistos en clase.
Pensando en ello, se presentan estas notas de clase que contienen ejercicios
rutinarios sobre integración para un curso de Cálculo Multivariado, en ellas se pretende trabajar básicamente tres elementos:
\begin{itemize}
  \item Integrar una función sobre un objeto: curva, superficie o sólido, parametrizando la región de integración.
    Esta idea es diferenciadora con respecto a la mayoría de textos de cálculo multivariado, pues la evaluación de integrales triples parametrizando el sólido es una estrategia muy poco utilizada, una ventaja es la obtención de límites constantes de integración.
  \item Se usan herramientas tecnológicas para abreviar el tiempo de ejecución, en particular, se ofrecen las instrucciones utilizadas con el asistente computacional Wolfram Mathematica para visualizar la región de integración o para la evaluación de integrales en diversas formas; dicha evaluación se hace incluso de manera numérica.
  \item Se ofrecen varias formas de resolver un mismo problema, en particular, se presenta la evaluación de integrales múltiples en diferentes órdenes.
\end{itemize}
Estos apuntes no pretenden, desde ningún punto de vista convertirse en un texto guía para un curso de Cálculo Multivariado.
Se desea hacer una compilación de problemas que ilustren su solución y que puedan servir como un material complementario de estudio.
La parte inicial de este material incluye los elementos teóricos básicos sobre los cuales se desarrollan los problemas.
Se agrega un capítulo con los comandos elementales para el uso del programa Wolfram Mathematica
y en los problemas presentados posteriormente se muestra las sintaxis requerida para utilizarlo.
En el tercer capítulo se introducen algunas parametrizaciones sencillas que se aplicarán posteriormente en los problemas que se
presentan en los capítulos 4, 5 y 6.
La integración de campos vectoriales y los teoremas relacionados se presenta en el capítulo 7.

Cada problema ofrece la figura asociada en diversas vistas y el cálculo que se requiere, se muestra de varias maneras.
Algunas de estas alternativas se ven muy engorrosas, pero es el resultado de la experiencia adquirida a traves de los años
en que los autores han orientado este curso.

El cambio de sistema de coordenadas para evaluar integrales triples, se hace tradicionalmente utilizando coordenadas esféricas y cilíndricas; no obstante, se incluyen otros sistemas de coordenadas en los anexos y una breve referencia en los ejercicios propuestos para motivar al lector al estudio de estos tópicos.

Este documento se editó en \LaTeX, sus más de 700 figuras se elaboraron en \verb"Mathematica" o con ayuda del paquete \verb"PSTricks" y su versión evolucionada \verb"psplotThreeD". La documentación puede consultarse en \textcite{vignault_pst-solides3d_2023}, \textcite{van_zandt_pst-plot_2016} y \textcite{vos_pst-func_2024}.

\begin{flushright}
  \smallcaps{Los autores}
\end{flushright}

\newpage